% Generated draft: requires mathematical and source review.
\documentclass[11pt]{article}
\usepackage[margin=25mm]{geometry}
\usepackage{fontspec}
\setmainfont{Latin Modern Roman}
\usepackage{amsmath,amssymb,mathtools,unicode-math}
\setmathfont{Latin Modern Math}
\usepackage{enumitem,needspace,xurl,hyperref}
\setlength{\parindent}{0pt}
\setlength{\parskip}{5pt}
\setlength{\emergencystretch}{3em}
\newcommand{\N}{\mathbb{N}}
\newcommand{\Z}{\mathbb{Z}}
\newcommand{\Q}{\mathbb{Q}}
\newcommand{\R}{\mathbb{R}}
\newcommand{\C}{\mathbb{C}}
\newcommand{\F}{\mathbb{F}}
\newcommand{\A}{\mathbb{A}}
\newcommand{\PP}{\mathbb{P}}
\newcommand{\E}{\mathbb{E}}
\newcommand{\eps}{\varepsilon}
\newcommand{\dd}{\,\mathrm{d}}
\begin{document}
\section*{UnsolvedMath ID: 30004351}
\textbf{Source problem code:} catalogue entry
\subsection{Definitions and mathematical statement}
Let \(G\) be a reductive algebraic group over \(\Q\), and let \(X_G\) denote the associated symmetric space referred to in the source. Assume that \(X_G\) has no complex structure, in the sense intended by the source.

\emph{Problem.} For every such \(G\), can one define a notion of modular forms for \(G\) admitting a Fourier expansion analogous to the Fourier expansion of Siegel modular forms?

\emph{Unresolved formulation qualifications.} The source does not specify the construction of \(X_G\), whether absence of complex structure means absence of a \(G(\R)\)-invariant complex structure, or the required properties of the proposed modular forms and their Fourier expansions. In particular, no arithmetic subgroup, level, weight or representation, regularity condition, growth condition, expansion domain, indexing set, coefficient space, or convergence requirement is prescribed. The question therefore requests the development of a theory; the supplied record does not determine a single precise existence assertion with specified axioms. These missing choices are retained as unresolved rather than supplied as additional hypotheses.
\subsection{Short English statement}
For a reductive group over \(\Q\) whose associated symmetric space has no complex structure, seek a theory of modular forms with Fourier expansions resembling those of Siegel modular forms. The intended meaning of complex structure and the required features of the theory and expansion remain unspecified.
\subsection{Sources}
\begin{itemize}
\item Local UnsolvedMath registry: supplied complete source record, permanent ID \(30004351\), source problem code \texttt{catalogue entry}, title \emph{Siegel Type Fourier Expansions Without Complex Structures}. The record belongs to the Oberwolfach Reports Open Problems collection and lists proposed year 2019; no individual proposer is supplied.
\item Supplied report citation: \emph{Modular Forms (2020)}, Oberwolfach Reports, DOI 10.4171/owr/2019/57, pp.\ 3529--3606. \url{https://doi.org/10.4171/owr/2019/57}. The supplied background describes these reports as documenting mathematical workshops and their open-problem discussions.
\item Supplied literature reference: Gourevitch--Gustafsson--Kleinschmidt--Persson--Sahi, \emph{Fourier coefficients}. \url{https://arxiv.org/abs/1811.05966}.
\item Supplied literature reference labeled ``OWR report'': \url{https://doi.org/10.4171/owr/2019/57}.
\item Literature assessment dated exactly \texttt{2026-08-22}: the supplied curation labels the problem partially solved and states that general automorphic-form and Fourier-coefficient theories for reductive groups provide broad analogues, including reduction principles for Fourier coefficients, but that no single holomorphic modular-form theory with the requested Siegel-like expansion exists for all such groups. This is supplied dated literature triage derived from the OWR curation, not an independently verified assessment; its time-sensitive status and cited claims require independent verification.
\end{itemize}
\end{document}
